Na het weesgegroetje
kwam uit moeders snoetje

een profetisch citaat.
Vader werd dan weer kwaad,

maar was er meestal niet,
als haar bede geschiedt.

“Wij”, zei de verteller,
“hebben een voorspeller

tussen onze oren.
Die kan een mens storen,

zelfs bij kalmte en rust.
Terwijl hij zelden sust.

Want klopt de voorspelling,
dan mist de verrassing.

Komt ie echter niet uit,
dan wordt dat fout geduid.

Namelijk als unfair.
Ons brein is de expert,

zo meent ie eigenwijs,
aangaande onze reis.

Ons supergeheugen
weet die pijn te heugen

en leert dat verwachting
leidt tot teleurstelling.

Dus worden we heel bang.
En dat is dan weer wrang,

omdat dit tekort doen
we steeds onszelf aandoen.

Daarbij is het “we” leeg.
Dat wat niet zijn hoop kreeg,

was slechts een verlangen.
Waar niets aan kan hangen

van wezenlijke aard.
Wat is het leven waard

als we moeten lijden?
Kan men zich bevrijden

uit de wrange strikken
van het vooruitblikken?

Boeddha bood een vluchtpad.
Jezus verbreedde dat

en verdroeg het lijden,
door het niet te mijden,

maar te transformeren.
De heer kon ons leren,

dat liefde ons verlost.
Lijden is het compost

voor diepe compassie.
De profeet zijn missie

was Jezus te volgen.
Zijn weg te vervolgen,

door de geest te bezien,
die het staag vooruitzien

als kenmerk in zich heeft.
Die zich voor jou uitgeeft,

maar slechts een klein deel is.
De profeet had gewis

zijn geest diepgaand verkend
en ons erin gekend.”

Emma’s naam werd verspreid
door haar losbandigheid.

Ze werd er om geliefd.
Een jongen werd verliefd.

Hij dacht als klasgenoot
dat ze haar hand aanbood,

omdat ze tegen hem
sprak met plechtige stem

en zelfs in het Latijn
tussen een raamkozijn

van Emma’s klaslokaal.
Ze zei mede brutaal

“Laten wij liefhebben!”
“Maar die spinnenwebben

spint ze om elke man,
met wie ze seksen kan,

sluw te hebben bediend”,
reageerde een vriend,

uit de kring van mannen,
die zijn klas bemanden

en rondom hem ontstond.
Dat snoerde niet zijn mond.

De amant zei snedig
“‘Ik bemin je eeuwig!’

Zegt ze toch niet zomaar.”
Zijn vriend zei: "Dat is waar!”

“Blijf in onze liefde”,
noemde de verliefde

ook nog als een bewijs
voor zijn verkregen prijs.

Zo noemde hij steeds meer
bewijzen van zijn eer.

“Ik zegen ons huw’lijk!”
gaf hem een gruwelijk

lot van ontnuchtering
door een medeleerling.
“Sprak ze tijdens Latijn?
Want haar uitspraken zijn

standaard oefenzinnen!”
De waarheid kwam binnen.

Jean-Paul Sartre's idee
van ‘de Ander’ deed mee.

Emmy bezag het recht
van een vrouw in de echt.

Zij zag de hetero-
normativiteit zo

duidelijk uitgebeeld
in wat had afgespeeld.

Derderangs aan de man,
lustte pap er pap van.

Al had mams nimmer zin,
ze bond meestal wel in.

Liggend in haar bedstee
kreeg Emmy alles mee.

De deur van haar slaapkast
zat immers zelden vast.

De slaapkamer van hout
was handig ingebouwd,

goedkoop en handzamer,
pal in de woonkamer.

De deur, vaak op een kier,
ontsloot steeds het vertier,

van haar oudelieden,
die ze ging bespieden.

Na haar moeders klagen
hoorde ze pa plagen

met: “dat is ‘s levens kelk”.
Tepeltje taptemelk

ging dan maar overstag
en staakte haar beklag.

Voor vaders hard bestaan
hield daags flink van bil gaan

hem draaglijk op de been,
waar even smart verdween.

Het bed van het tweetal
was zo hoog, kort en smal

als een relaxfauteuil.
Want voor haar ouwelui

stond slapen door liggen
gelijk aan doodliggen.

Men sliep dus halfzittend.
Tegen kussens liggend.

Dat vond men veel te krap
voor hechte gemeenschap.

Men vond de woonkamer
allicht aangenamer.

Pa nam dan ma daar mee.
Een krakende bedstee

en een sluipend gepiep,
nadat iedereen sliep,

werd zo een verrader.
Soms wilde haar vader

vaker worden verwend.
Een standaard argument
was haar bijgebleven:
“Maar ik ben zo even

niet fijn klaargekomen”
Het moest nogmaals stromen.

Vader kon niet lezen.
Hij moest bij mam wezen

om het voor te spellen
of het te vertellen

als een zaak op schrift stond.
Hij leefde bij de mond

en overleefde puur
op zijn starre natuur.

Struinend van dag naar dag
in non-stop zelfbeklag.

Maurice Merleau-Ponty’s
fenomenologie

berichtte dat ons lichaam
als instrument werkzaam

is bij ons bewustzijn
en in-de-wereld-zijn.

Door het voors te pellen
kon moeder voorspellen

wat ze precies moest doen
om de klaagkampioen

steeds in toom te houden
binnen haar huishouden.

's Ochtends was zijn pik
dik als een vaste prik.

Daar wist mam wel raad mee,
knielend naast de bedstee,

als ze naar hem afzonk
en zijn lading opdronk,

na wat wilde kussen.
Zo zagen de zussen

in hun tienerjaren
hun ouders vaak paren.

Destijds was dat gewoon
en geen onkies vertoon.

Emmy vond het smerig.
Emma werd begerig.

Ook normaal was naaktheid
in die vergane tijd,

als het functioneel was.
Men had geen ochtendjas,

maar slechts een nachtmuts aan
als men uit bed moest gaan

binnen het eigen huis.
Men vond dat niet onkuis.

In die situatie
vergeleek vaak Emmy

de fysieke franje
van de feloranje

stekels van de cactus
in het kruis van haar zus

met de zwarte krullen,
die de gleuf verhullen

van haar naakte moeder,
waar vader zijn poeder

dagelijks in dumpte
of zijn tong in wurmde.

Hoe zou zij zo’n leven
in lust gaan beleven?
